\origpage{137}{GRAPH MINORS. VI}
Moreover we can find all these homotopy classes efficiently, and so we have a polynomially bounded algorithm for DISJOINT CONNECTING SUBGRAPHS on the cylinder.

One can do much better than this. The algorithm above was given for its simplicity, but we sketch another approach. Once again we have at most
\begin{equation*}
\left| \bigcup \Delta \cap C_1 \right| \cdot \left| \bigcup \Delta \cap C_2 \right|
\end{equation*}
possibilities for the homotopy classes of the $H(s, t)$'s where $s, t$ are in the same cuff of $\Sigma$. Having chosen one of the possibilities to check, we can now apply reduction arguments like those in section 4 to reduce to the case where each $\delta \in \Delta$ has
\begin{equation*}
|\delta \cap C_1| = |\delta \cap C_2| = 1.
\end{equation*}
This case can be solved easily with at most $|\bigcup \Delta \cap C_1|$ applications of (5.7).

So far we have considered our problem DISJOINT CONNECTING SUBGRAPHS on a disc and on a cylinder. What about other surfaces? As mentioned in the introduction, it will be shown in [3] that there is a polynomially bounded algorithm to solve this problem on any fixed surface provided that $|\bigcup \Delta|$ is kept bounded. This is not a complete generalization of the results of this paper, because

\hspace*{12pt}(i) in this paper we do not need to keep $|\bigcup \Delta|$ bounded, and

\hspace*{12pt}(ii) we have no structural characterization in general.

However, a full analogue does exist when $\Sigma$ is the M\"obius band. Working with ROUTED FOREST CONTAINMENT, one reduces (using the methods of section 4) to the case when each component of the forest $H$ is a path with only its end-vertices in $bd(\Sigma)$, and none of these paths is homotopic to a portion of $bd(\Sigma)$. In that case, a necessary and sufficient condition for $G$ to contain a subgraph homotopic to $H$ is that there is no closed sequence
\begin{equation*}
v_0, R_1, v_1, R_2, \dots, v_{k-1}, R_k, v_k = v_0
\end{equation*}
of vertices and regions of $G$ (in which consecutive members are incident) which is ``homotopic'' to $bd(\Sigma)$ and with
\begin{equation*}
k < |bd(\Sigma) \cap V(G)|.
\end{equation*}
It is possible to base a polynomially bounded algorithm on that fact, but we omit the details.

Let $G$ be a planar graph, and let $s_1, t_1, \dots, s_k, t_k$ be distinct vertices of $G$,
